\section{Related Work}

{\noindent \bf Speech and Audio Synthesis.} Recent advancements in neural audio generation enabled computers to efficiently generate natural sounding audio. 
The first convincing results were achieved by autoregressive models such as WaveNet~\citep{oord2016wavenet}, at the cost of slow inference.
While many other approaches were explored~\citep{parallelwavegan,wavernn,goel2022s}, the most relevant ones here are those based on Generative Adversarial Networks (GAN)~\citep{melgan, parallelwavegan, hifigan, hifi++} were able to match the quality of autoregressive  by combining various adversarial networks operate at different multi-scale and multi-period resolutions. Our work uses and extends similar adversarial losses to limit artifacts during audio generation.

% and SampleRNN~\cite{mehri2016samplernn} generate high-quality audio in the waveform domain, one sample at a time, resulting in slow inference. 
% WaveRNN~\cite{wavernn} enabled fast synthesis by training a compact neural network further optimized via sparsification. Another line of auto-regressive model involves in~\cite{goel2022s}, in the authors propose an audio generative model based on the Structured State Space Sequence model~\cite{gu2021efficiently}. Feed-forward networks were additionally suggested to further speed up inference, via knowledge distillation from an autoregressive teacher-model~\cite{parallelwavenet,ping2018clarinet, kim19flowavenet}. Two other lines of related work concerns with Flow base neural vocoders~\cite{waveglow} and diffusion based neural vocoders~\cite{wavegrad, diffwave}. Recently, generative adversarial networks (GANs)~\cite{melgan, parallelwavegan, hifigan, hifi++} were able to match the quality of autoregressive and large feed-forward models by combining various adversarial networks operate at different multi-scale and multi-period resolutions. 

{\noindent \bf Audio Codec.} Low bitrate parametric speech and audio codecs have long been studied~\citep{atal1971speech, juang1982multiple}, but their quality has been severely limited. Despite some advances~\citep{griffin1985new, mccree19962}, modeling the excitation signal has remained a challenging task. The current state-of-the-art traditional audio codecs are Opus~\citep{valin2012definition} and Enhanced Voice Service (EVS)~\citep{dietz2015overview}. These methods produce high coding efficiency for general audio while supporting various bitrates, sampling rates, and real-time compression. 

\looseness=-1
Neural based audio codecs have been recently proposed and demonstrated promising results~\citep{kleijn2018wavenet,valin2019real, lim2020robust, kleijn2021generative, zeghidour2021soundstream, omran2022disentangling, lin2022speech, jayashankar2022architecture, livariable, jiang2022end}, where most methods are based on quantizing the latent space before feeding it to the decoder. In~\citet{valin2019real}, an LPCNet~\citep{valin2019lpcnet} vocoder was conditioned on hand-crafted features and a uniform quantizer. \citet{garbacea2019low} conditioned a WaveNet based model on discrete units obtained from a VQ-VAE~\citep{van2017neural, razavi2019generating} model, while \citet{skoglund2019improving} tried feeding the Opus codec~\citep{valin2012definition} to a WaveNet to further improve its perceptual quality. \citet{jayashankar2022architecture, jiang2022end} propose an auto-encoder with a vector quantization layer applied over the latent representation and minimizing the reconstruction loss, while \citet{livariable} suggested using Gumbel-Softmax (GS)~\citep{jang2016categorical} for representation quantization. The most relevant related work to ours is the SoundStream model~\citep{zeghidour2021soundstream}, in which the authors propose a fully convolutional encoder decoder architecture with a Residual Vector Quantization (RVQ)~\citep{gray1984vector, vasuki2006review} layers. The model was optimized using both reconstruction loss and adversarial perceptual losses.

{\noindent \bf Audio Discretization.}
Representing audio and speech using discrete values was proposed to various tasks recently. \citet{dieleman2018challenge, dhariwal2020jukebox} proposed a hierarchical VQ-VAE based model for learning discrete representation of raw audio, next combined with an auto-regressive model, demonstrating the ability to generate high quality music. Similarly, \citet{lakhotia2021generative, kharitonov2021text} demonstrated that self-supervised learning methods for speech (e.g., HuBERT~\citep{hsu2021hubert}), can be quantized and used for conditional and unconditional speech generation. Similar methods were applied to speech resynthesis~\citep{polyak2021speech}, speech emotion conversion~\citep{kreuk2021textless}, spoken dialog system~\citep{nguyen2022generative}, and speech-to-speech translation~\citep{lee2021direct,lee2021textless,popuri2022enhanced}.